% Proof-audit section for the recovered prefix250 construction.
% Requires amsmath, amssymb, and hyperref in the enclosing document.
% Audited against ProveIt 80e26b89e905bf3e5dfc322befdc1c6eaf25a109.
\section{Whole-frame simulation and the ZFC equivalence}
\label{sec:frame-calibration}

This section supplies the unbounded composition argument for the retimed
250-state table.  It uses the three new register-operation identities
proved in the local-contract section, the unchanged register-marking
identity, and the exhaustive symbolic instruction certificate.  Finite
literal runs provide an independent check of the identities and of the
bootstrap.  They do not replace the induction arguments.

All numerical state labels in this section are original labels.  The
machine starts at state~97 on the blank tape; its recurrent dispatch
state is~96.  Its original-label transition file has SHA-256
\begin{quote}\small\ttfamily
fd455304bad82a6da7ab83149352d0b1db27557582c4cc6312de003c96ff4e43.
\end{quote}
The 250-state count includes the initializer.  The construction and
the preceding shift/pulse investigation are credited separately from
this independent proof audit.

\subsection{The complete recurrent invariant}

Let $R:\mathbb N\to\mathbb N$ be the logical register file, in the order
$(N,T,W,A,D,H,X,Y,L)$ on indices $0,\ldots,8$, with
$R(i)=0$ for $i\ge9$.  The physical primary list is
\[
 V(R)=(R(2),R(8),R(1),R(4),R(3),R(7),R(0),R(5),R(6)).
\]
Thus the physical positions of the nine logical registers are
$(6,2,0,4,3,7,8,5,1)$, exactly as in permutation252.  Write
$B(v)=0\,1^{v+1}$ and extend $B$ by concatenation to finite lists.

A recurrent frame consists of an integer head position $h$, a finite
program-counter word $P$, a natural header distance $d$, the primary
registers $R$, and a \emph{nonempty finite} spare list $S$.  Its complete
infinite tape is
\[
 0^{\mathbb Z_{<h-1}}\;1\;P\;0^{d-|P|}\;1\;B(V(R)S)\;0^\omega,
 \qquad d\ge40,\quad |P|\le d.
\]
The left marker is at $h-1$, the header is at $h+d$, and the first
register separator is at $h+d+1$.  The head is at $h$ and the state
is~96.  This notation specifies every tape cell; it does not allow
unexamined nonzero data beyond a finite window.  Trailing zeros in a
finite representation of either exterior are immaterial.

The lower bound concerns the \emph{total distance from the head to the
header}.  In the original Lean definition
\code{Permutation252Frame.DispatchFrame}, the equivalent bound is
$40\le |P|+\mathit{gap}$, not $40\le\mathit{gap}$.  At an empty
root, these numbers coincide.  The value~40 is the sufficient bound
used by the complete instruction certificate; no optimality claim
is made for that margin.

There is one additional, essential program-counter invariant.  Extend
$P$ by zeros on the right and follow the literal register decoder to a
leaf.  If that leaf consumes $\ell$ bits, every bit of $P$ at an index
$\ge\ell$ is zero.  This is the predicate
\code{RegisterPathClean.Clean}.  It holds for the empty word, and
\code{RegisterPathClean.step\_clean} proves that every successful
logical step produces it.  In particular, an arbitrary bit string
that happens to fit physically before the header is not automatically
a valid recurrent source frame.

The physical spare suffix and the zero logical tail have different
meanings.  The decoder selects only one of the first nine physical
blocks.  A nonzero spare is not the value of a logical register with
index at least~9.  The encoding $R(i)=0$ beyond~8 is maintained by
logical updates to indices below~9; the physical suffix $S$ is retained
as tape data with its own invariant.

\subsection{Why every legal dispatch is covered by a leaf certificate}

The inherited literal program has 6,459 binary decoder nodes and 6,460
leaves.  Every leaf path has length at most~34.  The new symbolic
certificate covers every leaf, and both guards at every decrement
leaf, for 8,391 cases in total.  Tree completeness must be verified
as well as the case count: an omitted decoder branch cannot be
silently treated as a halt.  The inherited interfaces
\code{SiteCoverage.allSites\_walk\_total} and
\code{SiteCoverage.allSites\_walk} express precisely these two
requirements.

For a logical word $P$ and its consumed depth $\ell$, put
\[
 Q=\operatorname{take}_{\ell}(P0^\omega),\qquad n=d-40.
\]
The clean-path invariant and $\ell\le34\le40\le d$ give the literal
finite-word identity
\[
 P\,0^{d-|P|}=Q\,0^{d-\ell}
                =Q\,0^{n+40-|Q|}.
\tag{1}\label{eq:clean-rebase250}
\]
If $|P|<\ell$, the missing consumed bits come from the existing gap.
If $|P|>\ell$, the removed bits are zero by cleanliness.  This proves
that the complete binary source tape is exactly the source tape of the
appropriate leaf certificate.  It is the ordinary word argument
formalized in \code{RegisterPathClean.gap\_consume} and
\code{Permutation252Simulation.source\_represents}; the new table
does not change that argument.

There is a second harmless representation issue at the output.  The
Python certificate writes an incremented path only through its last
remaining one, whereas \code{RegisterPath.incPath} retains the
erased trailing zeros.  Both encode the same zero-padded word.  For
example, $0101$ and $010100$ have the same bit at every position after
zero extension.  Moving the omitted zeros into the gap leaves the
complete tape unchanged.  On carry overflow the logical successor is
an all-zero word of its old length; its physical empty-word
representation has the same zero-padded bits.  Equation~\eqref{eq:clean-rebase250}
therefore handles both source and target representations without
changing the register program.

\subsection{The symbolic certificate is uniform, not a numerical sample}

The symbolic tape contains literal bits, uniform runs of lengths
$an+c$ with $a,c\ge0$, and one register-file object.  Put $U=V(R)$.
The latter has
exactly one of the states: unchanged, marked at a specified physical
register, incremented there, or decremented there.  Its meanings in
this construction are
\[
\begin{array}{c|l}
\text{effect}&\text{complete represented register file}\\\hline
\text{unchanged}&B(US)\\
\text{increment at physical }r&B((U[r\mapsto U_r+1])\,(S\concat[1]))\\
\text{positive decrement at physical }r&B((U[r\mapsto U_r-1])\,S).
\end{array}
\]
Here a logical instruction selecting $i$ marks physical position
$r=\pi(i)$ for the permutation stated above.  The positive decrement
has the additional one-cell movement of the
register-file boundary established by its local identity.

Each rule has an elementary interpretation on complete tapes.
A literal rule is the specified transition.  Peeling a run uses
$an+c=1+(an+c-1)$ only when $c>0$, so it is valid also when $n=0$.
A scan rule is induction on the actual length of a uniform run, and
is permitted only for the indicated self-loop and direction.
Recentring the zipper, deleting a provably empty run, and padding
the blank exterior change the representation, not the configuration.
Adjacent uniform segments are merged by addition of their lengths.
The register-file rules are exactly the four stated marking and
operation contracts; a decrement rule additionally checks that its
guard names the same physical register as the marked block.

Consequently induction on an accepted symbolic execution gives, for
\emph{every} $n\ge0$, every natural-valued primary register vector
satisfying its named guard, and every nonempty finite spare list, a
finite literal run with exactly the stated complete output frame.
The single-effect representation prevents a register update from being
silently reapplied.  An unrecognized boundary case or a run whose
nonzero length cannot be justified is rejected by the checker.

The clock supplied by a scan can be zero, so positivity is checked
separately.  Every accepted instruction/guard certificate contains
at least 13 literal transitions.  All other rules denote
nonnegative-length computations.  Thus every nonhalting logical
step has a strictly positive finite binary duration, for all
parameter values.  This claim follows from the rule interpretation
and the finite certificate, not from the sizes used in regression
tests.  The Python checker is not itself a Lean or Rocq kernel
checker; its mathematical rule-soundness proof is part of the
ordinary proof supplied here.

\subsection{Preservation of distance, spares, and primary values}

Let $\delta$ be~1 exactly for a positive decrement and~0 otherwise.
Let $\varepsilon$ be~1 exactly when the program-counter update carries
past the left marker, and~0 otherwise.  The complete instruction
certificate establishes the endpoint coordinates
\[
 h'=h-\varepsilon,\qquad
 d'=d+\delta+\varepsilon.
\tag{2}\label{eq:distance-update250}
\]
Thus the absolute header coordinate changes by $\delta$:
$h'+d'=h+d+\delta$.  A positive decrement moves the register file
right by one; an overflow moves the left marker and dispatch head
left by one.  The two changes are accounted for separately.

For an increment, the selected primary value increases by one and
$S'=S\concat[1]$.  For a positive decrement the selected primary value
decreases by one and $S'=S$.  For a zero decrement, break, or
continue, the primary vector and spare list are unchanged.  The
program-counter words are precisely the inherited Wade semantics:
increment advances its consumed path; positive decrement does the
same; zero decrement performs a one-level break; break strips
through the designated number of zero bits and then advances;
continue strips through the designated number of one bits.

Equations~\eqref{eq:clean-rebase250} and
\eqref{eq:distance-update250} restore a recurrent frame for the
logical successor.  The distance stays at least~40.  The spare list
stays nonempty, and actually grows only on increment.  Its values
need not be bounded or zero.  At every decrement of the last
primary block, a following block is therefore present.  This last
condition is necessary for the zero-decrement rollback probe: with
no following block it can take the end-of-file branch instead.
The new invariant is thus slightly stronger than the arbitrary,
possibly empty spare invariant of permutation252, and is proved
inductive here.  No finite reserve can be exhausted by a long run.

Logical updates preserve $R(i)=0$ for all $i\ge9$, because every
selected logical register lies below~9.  The permutation and its
inverse ensure that the actual primary block modified is the
intended logical register.  These finite index identities are
already expressed by \code{Permutation252Frame.nth\_vregs\_append},
\code{update\_vregs\_append}, and \code{tail\_update}.

The relevant boundary is a \emph{complete dispatch frame}, not merely
the numerical state~96.  A local decoder execution may revisit
state~96 before it has restored that frame.  The certificate seeks
the whole target configuration, so these intermediate visits do
not produce false logical steps.

\subsection{Unbounded halting equivalence from positive segments}

Let $\mathcal R(a,c)$ mean that the binary configuration $c$ has the
complete recurrent frame for the clean logical register
configuration $a$.  The preceding results give four facts:
\begin{enumerate}
\item If $\mathcal R(a,c)$, then $c$ is active, in state~96.
\item If the logical step is $a\mapsto b$ and $\mathcal R(a,c)$,
there is $k>0$ such that $\mathcal R(b,M^k(c))$.
\item If the logical step has no successor, full decoder coverage
implies that its decoded leaf is a literal halt leaf.
\item From a related literal halt leaf, a finite binary segment
executes the designated halting transition.
\end{enumerate}
The final write and head movement of that transition are retained.
Afterward the semantics is absorbing.

The forward implication is finite induction: compose the positive
segments along a logical run reaching a halt leaf, then compose the
literal halting segment.  This yields a finite binary halt witness.

For the reverse implication, a particularly direct proof is strong
induction on a binary halting time $t$.  Suppose
$\mathcal R(a,c)$ and $M^t(c)$ is halted.  If the logical step has no
successor, fact~3 supplies the logical halt.  Otherwise choose the
segment of fact~2, of length $k>0$, reaching an active related
configuration $c'$.  Necessarily $k<t$: if $k\ge t$, absorption of
the halt would make $c'$ halted.  The remaining binary run from
$c'$ halts after $t-k<t$ steps.  Apply the induction hypothesis and
prepend the one logical step.  This establishes
\[
 \mathcal R(a,c)\ \Longrightarrow\
 \left[\begin{aligned}
 &\bigl(a\text{ reaches a logical halt leaf}\bigr)\\
 &\qquad\longleftrightarrow
 \bigl(c\text{ eventually halts under }M\bigr)
 \end{aligned}\right].
\tag{3}\label{eq:global250}
\]
No uniform clock bound, infinite choice of segments, or unproved
fairness principle is involved.  The generic argument is already
kernel checked in \code{RuntimePerm252.PositiveSimulation},
\code{binary\_to\_abstract}, and \code{halts\_iff}; its
instantiation for the new table uses the new contracts and frame.

\subsection{The verified blank prefix enters an admitted dirty root}

The literal blank run reaches its recorded complete frame at
585,431 transitions:
\[
 q=96,\quad h=-3,\quad d=228,\quad P=[],\qquad S=[1]^{210}.
\]
Its physical primaries are $(0,2,0,4,0,0,1,0,0)$, so
\[
 R_*=(N,T,W,A,D,H,X,Y,L)=(1,0,0,0,4,0,0,0,2).
\tag{4}\label{eq:dirty-root250}
\]
The marker is at $-4$, the header at~225, and there are 438 one-cells.
The published finite support lists all of them, with every other
cell blank.  The empty program counter is clean, $228\ge40$, and
the 210 spare blocks are nonempty, so every clause of
$\mathcal R(\langle[],R_*\rangle,c_*)$ holds.

The dirty-start theorem does not require the six numerical fields
$N,A,D,X,Y,L$ to vanish.  It requires the formula invariant
\[
 C_{254}(T)\ \wedge\ C_{254}(H)\ \wedge\operatorname{Good}(W).
\]
Here $T=H=0$ are the known truth code, and $W=0$ is the empty stack.
The exact witnesses are \code{Code254.C254.truth},
\code{Code254.C254.truth}, and \code{Enum254.good\_zero}.
The nonzero fields $N=1$, $D=4$, and $L=2$ therefore create no
unproved soundness obligation.

For clarity, the inherited root comparison is the composition
\[
\begin{split}
 \operatorname{Halts}_{\rm Prog254}([],R_*)
 &\longleftrightarrow
   \operatorname{Halts}_{\rm pathDag254}([],R_*)\\
 &\longleftrightarrow C_{254}(F),\qquad F=1.
\end{split}
\tag{5}\label{eq:root-calibration250}
\]
The first line is
\code{Bisim254.prog254\_halts\_iff R}, valid for every register
file $R$.  The second is
\code{DirtyProgram254.path\_halts\_iff}, applied to
\code{L2.enc} of the nine-field state in
\eqref{eq:dirty-root250}.  It is not an application of a theorem
restricted to the older zero root or to the public $N=2$ start.

The dirty theorem uses an explicit absorption map:
\code{DirtyRoot254.virtual} makes scratch and accumulator fields
clean while preserving the continuation after the first fetch and
read routines.  Its theorem
\code{DirtyRoot254.stage\_virtual} proves equality of those
continuations.  The formula fields are preserved by
\code{virtual\_formulas}.  Completeness from an arbitrary
candidate counter is proved by
\code{EnumSafeStart254.complete\_from}: a padded refutation
program can be chosen with code above the current counter, and
every later candidate is reached unless an earlier candidate
already halts.  Thus no possible refutation is lost by starting
with $N=1$ or with a nonzero list pointer.

Adding or removing an actual finite prefix preserves eventual
halting for an absorbing machine.  Combining the blank prefix,
\eqref{eq:global250}, and \eqref{eq:root-calibration250} gives
\[
 \operatorname{Halt}(M_{250},\mathrm{blank})
 \longleftrightarrow C_{254}(F).
\tag{6}\label{eq:binary-code250}
\]
The right side is a finite numerical derivability predicate.
\code{Code254Proof.c254\_iff\_checked} additionally represents it
as the existence of a natural proof code accepted by the exact
finite checker.  This is an inherited theorem, not a newly
assumed equivalence between an informal calculus and ZFC.

\subsection{Ordinary ZFC proof of the calibration}

Write $\operatorname{Ref}_T$ for the existence of a finite formal
refutation in a fixed recursive presentation of $T$, and
$\operatorname{Con}(T)=\neg\operatorname{Ref}_T$.  Fix the native
theory $T_{\rm nat}$ to be exactly
\code{NativeTheory254.NativeAx}: universal closures of CP,
every native-formula GIC instance, and raw irreflexivity.
The GIC scheme ranges over first-order formulas.  It is not an
axiom quantifying over every external binary relation.

The inherited mathematical calibration consists of
\[
 C_{254}(F)
 \longleftrightarrow \operatorname{Ref}_{T_{\rm nat}}
 \longleftrightarrow \operatorname{Ref}_{\rm ZF}
 \longleftrightarrow \operatorname{Ref}_{\rm ZFC}.
\tag{7}\label{eq:native-zfc250}
\]
The exact endpoints are
\code{NativeTheory254.c254\_refutation\_iff},
\code{WFTheory.native\_refutation\_iff\_zf} (or
\code{WFTheory.c254\_refutation\_iff\_zf} for the direct code
comparison), and
\code{ChoiceCalibration254.zf\_refutation\_iff\_zfc}.
The independently composed endpoint
\code{ConstructibleChoice.zf\_refutation\_iff\_zfc\_refutation}
is used by \code{BinaryZFC252}.

Here is the set-theoretic reading that makes the strength of
\eqref{eq:native-zfc250} precise.  Work \emph{inside ZFC} and code
syntax and proofs by finite sequences of natural numbers.  Each
theory is a set of first-order sentences in a countable language.
First-order soundness and the Henkin completeness construction
are theorems of ZFC for such coded theories.  In particular,
consistency supplies a nonempty set-coded model, which can be
taken countable; it supplies neither a transitive model nor an
externally well-founded model.

For the native-to-ZF comparison, the two model constructions used
by the inherited proof can be performed on those sets:
\begin{enumerate}
\item Every ZF model satisfies the native axioms.  Raw
irreflexivity follows from Foundation.  Suitable internal limit
rank segments supply the containers.  The internally derived
Collection scheme supplies the literal GIC instances.  These
are the contents of \code{ZFNative.native\_of\_zf}, based on
\code{ZFRank.containers} and \code{ZFRank.collection}.
Consequently a native refutation rules out all ZF models, and
completeness gives a ZF refutation.
\item From a native model, its CP/GIC reduct supplies the
bisimulation quotient used in \code{QuotientZF0}.  As a
quotient of a set this is again a set.  The quotient has the
required extensional ZF axioms without Foundation, together
with the actual power-set construction.  Its \emph{internally}
well-founded subdomain supplies Foundation and retains the
remaining axioms.  This is
\code{WFModel.quotient\_models\_ZF}; its nonemptiness is
\code{WellFoundedPart.quotient\_domain\_nonempty}.
Thus a ZF refutation rules out all native models, and
completeness gives a native refutation.
\end{enumerate}

The word ``internally'' is indispensable in item~2.  In the
quotient model $(D,E)$, the predicate is
\[
\begin{split}
 \operatorname{WF}_E(x)\;\Longleftrightarrow\;
 \exists t\in D\;[\;&xEt\ \wedge\ t\text{ is }E\text{-transitive}\\
 &\wedge\ \forall s\in D\,
   (s\subseteq_E t\ \wedge\ s\ne_E\varnothing
      \Rightarrow s\text{ has an }E\text{-minimal member})].
\end{split}
\]
This is one first-order formula, explicitly defined as
\code{WellFoundedPart.wfForm}; the subset quantifier ranges
over elements represented in $D$.  It is not external
well-foundedness of $E$.  Separation in the metatheory forms
$\{x\in D:\operatorname{WF}_E(x)\}$.  The specific GIC witness
instance proves closure under sets of such elements.  The
relativized Separation and Replacement schemes are handled by
explicit formula translations in \code{WFRelativization}
and \code{WFModel.domain\_replacement}.  Infinity is
obtained by separating this predicate from an internal
inductive set.  No assumption of full Collection in an
arbitrary ZF-without-Foundation model is being substituted.

For ZF versus ZFC, inclusion of the ZF axioms in ZFC gives one
refutation direction.  For the other, assume a ZFC refutation
and suppose ZF were consistent.  Henkin completeness gives
a set model $\mathcal A$ of ZF.  Form its internal constructible
subdomain
\[
 L^{\mathcal A}=\{a\in |\mathcal A|:
                       \mathcal A\models a\in L\}.
\]
It is a subset of the set $|\mathcal A|$, even though
$\mathcal A$ regards its own $L$ as a class.  The usual
constructible-hierarchy proofs establish, uniformly for
first-order formula codes, that this induced structure
satisfies ZF and Choice.  In the pinned development these are
\code{ConstructibleZF.L\_models\_ZF},
\code{ConstructibleAbs.vEqL}, and
\code{ConstructibleOrder.choice\_of\_VL}, composed in
\code{ChoiceCalibration254.constructible\_models\_ZFC}.
There is an explicit inhabitant, supplied by
\code{constructible\_nonempty}.  Soundness contradicts the
assumed ZFC refutation.  Classical logic therefore gives a ZF
refutation.

All structures in this argument are set coded.  Their
satisfaction predicates are defined by recursion on finite
formula syntax; this is satisfaction for \emph{set structures},
not a truth predicate for the entire universe.  Quotients,
definable subdomains, and finite formula translations are
ordinary ZFC constructions.  Uniform induction on formula
codes handles the axiom schemes.  The internally ill-founded
or nonstandard nature of a Henkin model creates no extra
premise: all the interpreted predicates are evaluated in
that model and all formula inductions are performed in the
metatheory.  This gives an ordinary ZFC proof of
\eqref{eq:native-zfc250}, with the inherited interpretation
lemmas as its identified mathematical ingredients.

The machine part is simpler to internalize at this strength.
Finite binary tapes, register files, paths, run segments, and
proof trees are sets coded by finite sequences of natural
numbers.  The transition function and the finite trace
predicate are primitive-recursive operations.  The local
word identities use induction on finite block lists; the
symbolic rule proof uses induction on a finite execution;
the global reverse direction uses induction on a natural
halting time.  ZFC proves these inductions and the required
recursions.  The specific bootstrap and all-sites certificate
are finite computational equalities for the displayed
tables.  Thus equations~\eqref{eq:binary-code250} and
\eqref{eq:native-zfc250} compose in ZFC to give
\[
 \boxed{\quad
 \mathrm{ZFC}\ \vdash\
 \left[\neg\operatorname{Halt}(M_{250},\mathrm{blank})
 \ \longleftrightarrow\ \operatorname{Con}(\mathrm{ZFC})\right].
 \quad}
\tag{8}\label{eq:target250}
\]
Neither direction assumes that ZFC is consistent.  The proof
establishes an equivalence, not either side of it in isolation.

\subsection{Proof status and reproducibility boundary}

The claim in \eqref{eq:target250} is an ordinary mathematical
proof with explicit finite computational certificates and
inherited interpretation lemmas.  The inherited permutation252
endpoint has an external Lean proof according to its pinned
repository verification records.  This audit does not claim
to have rebuilt that dependency closure.  The new 250-state
endpoint has not been delivered as a closed Lean or Rocq
theorem, nor as a serialized first-order ZFC derivation.

In particular, a Lean axiom report containing only
\code{propext}, \code{Classical.choice}, and
\code{Quot.sound} is not an automatic translation of an
arbitrary universe-polymorphic Lean proof into ZFC.  The
ZFC claim above rests on the specific set-coded mathematical
argument, with countable completeness and set-sized model
constructions, not on such a general translation.  The
repository's stronger internal PA proof-code and finite-trace
formalization goal remains a separate task.

Under the additional metatheoretic hypothesis
$\operatorname{Con}(\mathrm{ZFC})$, the machine does not
halt, and G\"odel's second incompleteness theorem prevents
ZFC from proving its nonhalting.  To infer full two-sided
independence by also excluding a ZFC proof of the false
existential halting assertion, one can assume
$\Sigma_1$-soundness.  Consistency alone does not generally
exclude a theory proving its own inconsistency.  These
independence qualifications do not change
\eqref{eq:target250}.

\subsection{Pinned source interfaces}

All ProveIt paths in the following list are relative to
\code{Computability/BusyBeaver/Lean/BusyBeaver/ZFC/}, at
commit
\href{https://github.com/VladimirReshetnikov/ProveIt/tree/80e26b89e905bf3e5dfc322befdc1c6eaf25a109/Computability/BusyBeaver/Lean/BusyBeaver/ZFC}{\code{80e26b89e905bf3e5dfc322befdc1c6eaf25a109}}.
The exact table and source hashes should accompany the
machine package so that later repository changes do not
silently alter the dependencies.
\begin{itemize}
\item \code{RegisterPath.lean},
\code{RegisterPathClean.lean},
\code{SiteCoverage.lean}: literal path semantics,
clean-path preservation, and complete leaf coverage.
\item \code{Permutation252Frame.lean},
\code{Permutation252Simulation.lean},
\code{RuntimePerm252.lean}: register placement,
complete-tape source rebasing, and positive-segment
halting reflection.  The frame is strengthened here by
requiring a nonempty spare suffix.
\item \code{Bisim254.lean},
\code{DirtyRoot254.lean},
\code{DirtyProgram254.lean},
\code{EnumSafeStart254.lean}: literal decoder to
219-node DAG, scratch absorption, safe dirty roots,
and completeness from arbitrary candidate counters.
\item \code{Code254Proof.lean},
\code{NativeTheory254.lean}: the finite numerical
proof checker and the exact CP/all-Form-GIC/IR theory.
\item \code{ZFNative.lean},
\code{WellFoundedPart.lean},
\code{WFRelativization.lean},
\code{WFModel.lean}, \code{WFTheory.lean}:
the two native/ZF model directions, with the internal
well-founded predicate and the literal scheme translations.
\item \code{ConstructibleChoice.lean},
\code{ChoiceCalibration254.lean},
and \code{BinaryZFC252.lean}.
These supply constructible-subdomain
Choice and the two independently composed external
ZFC consistency comparisons for the inherited machine.
\end{itemize}
